\chapter*{Remerciements}
\addcontentsline{toc}{chapter}{Remerciements}
\markboth{Remerciements}{Remerciements}

Au terme de ce stage d'ingénieur d'immersion professionnelle au sein du \textbf{Groupe OCP}, je tiens à exprimer ma sincère et profonde gratitude à toutes les personnes qui ont contribué, par leur disponibilité, leur expertise et leurs précieux conseils, à la réussite de ce projet de Jumeau Numérique.

J'adresse tout d'abord mes vifs remerciements à la direction du \textbf{Groupe OCP S.A.} pour m'avoir accueilli dans un environnement industriel d'excellence propice à l'innovation technologique.

J'exprime ma profonde reconnaissance à mon encadrant professionnel, \textbf{[ENCADRANT PROFESSIONNEL]}, pour son accueil bienveillant, sa confiance constante, ses orientations stratégiques éclairées et son accompagnement technique tout au long de la conception et du développement de ce prototype. Ses explications sur les processus de traitement minier ont été déterminantes pour aligner la modélisation logicielle avec les exigences opérationnelles du terrain.

Mes sincères remerciements vont également à mon encadrant pédagogique à l'ENSA Safi, \textbf{[ENCADRANT PÉDAGOGIQUE]}, pour son suivi académique rigoureux, ses conseils méthodologiques pertinents et son soutien sans faille durant toutes les phases de ce travail.

Je remercie chaleureusement l'ensemble des équipes techniques, automaticiens et ingénieurs procédés du site d'accueil, dont l'esprit d'équipe, l'écoute et le partage d'expérience ont grandement enrichi mon immersion dans le milieu de l'Industrie 4.0.

Enfin, j'adresse ma vive gratitude au corps professoral et administratif de l'\textbf{École Nationale des Sciences Appliquées de Safi (ENSA Safi)} pour la qualité de la formation dispensée, qui m'a doté du socle scientifique, technique et méthodologique indispensable à l'accomplissement d'un tel projet d'ingénierie logicielle.